\section*{अध्याय \ref{ch:g11:subduction-zones} --- क्षेपण क्षेत्र}

\begin{solution}{exo:g11:subduction-zones:1}
\omterm{def:g11:subduction-zones:subduction}{क्षेपण} \omterm{def:g11:ocean-ridges:oceanic-lithosphere}{महासागरीय स्थलमंडल} की किसी प्लेट का किसी दूसरी प्लेट के नीचे, वापस
प्रावार में, डूबना है; \omterm{def:g11:subduction-zones:subduction}{क्षेपण क्षेत्र} वह अभिसारी प्लेट-सीमा है जहाँ यह
होता है, और इसे एक खाई चिह्नित करती है। \omterm{def:g11:ocean-ridges:oceanic-lithosphere}{महासागरीय स्थलमंडल} डूबता है: पुराना होने पर
वह ठंडा, मोटा और अपने नीचे के \omterm{def:g11:lithosphere-asthenosphere:lithosphere}{दुर्बलतामंडल} से अधिक घना होता है।
महाद्वीपीय भूपर्पटी डूबने के लिए बहुत हल्की है।
\end{solution}

\begin{solution}{exo:g11:subduction-zones:2}
(1) उदाहरण के लिए टोंगा--फ़िजी क्षेत्र, जापान से मारियाना तक का पश्चिमी
प्रशांत, इंडोनेशिया, और दक्षिण अमेरिका। (2) हमेशा ऊपरी प्लेट की
ओर, कभी महासागर वाली ओर नहीं। (3) उथले: मानचित्र पर मध्य-अटलांटिक कटक के
सभी भूकंप \qty{70}{km} से उथले हैं।
\end{solution}

\begin{solution}{exo:g11:subduction-zones:3}
$7031 \times 0.21 \approx 1480$ भूकंप \qty{70}{km} से अधिक गहरे;
$7031 \times 0.069 \approx 485$ \qty{300}{km} से अधिक गहरे।
\end{solution}

\begin{solution}{exo:g11:subduction-zones:4}
\omterm{def:g11:subduction-zones:eclogite}{नीलशिस्ट} \omterm{def:g11:ocean-ridges:oceanic-lithosphere}{महासागरीय भूपर्पटी} के बेसाल्ट से बनी कायांतरित चट्टान है, जो
अपेक्षाकृत ठंडी रहते हुए उच्च दाब तक दबाई गई हो, और जिसे ग्लॉकोफ़ेन नीला रंग देता है;
\omterm{def:g11:subduction-zones:eclogite}{एक्लोजाइट} वह घनी कायांतरित चट्टान है जो बेसाल्ट और \omterm{def:g11:ocean-ridges:gabbro}{गैब्रो} प्रावार की गहराइयों पर
बन जाते हैं, लाल गार्नेट और हरे पाइरॉक्सीन से बनी। दोनों
\omterm{def:g11:ocean-ridges:oceanic-lithosphere}{महासागरीय भूपर्पटी} से आती हैं और उसका रासायनिक संघटन बनाए रखती हैं; भिन्न हैं
खनिज, जो नए दाब और तापमान में ठोस अवस्था में बने,
और घनत्व।
\end{solution}

\begin{solution}{exo:g11:subduction-zones:5}
महासागरीय तली: खाई की ओर चलता पुराना \omterm{def:g11:lithosphere-asthenosphere:lithosphere}{स्थलमंडल}। खाई: प्लेट नीचे मुड़ती है
और डूबना शुरू करती है। अग्र-चाप: डूबती प्लेट से खुरचे गए अवसाद ढेर
होते हैं; नीचे, दोनों प्लेटें प्लेट-सीमा पर एक-दूसरे के पास से सरकती हैं,
जहाँ महान भूकंप होते हैं। ज्वालामुखीय चाप: \omterm{def:g11:subduction-zones:slab}{स्लैब} के ऊपर ज्वालामुखियों की
एक पंक्ति। पश्च-चाप: कुछ चापों के पीछे ऊपरी प्लेट खिंचती है।
\end{solution}

\begin{solution}{exo:g11:subduction-zones:6}
$\tan\delta = (508 - 135)/(800 - 300) = 373/500 \approx 0.746$, इसलिए
$\delta \approx \ang{37}$।
\end{solution}

\begin{solution}{exo:g11:subduction-zones:7}
(1) $\tan\ang{45} = 1$: \qty{300}{km} गहरी। (2) खाई से \qty{670}{km}
पर। (3) \omterm{def:g11:subduction-zones:slab}{स्लैब} की $670\sqrt{2} \approx \qty{948}{km}$ लंबाई।
\end{solution}

\begin{solution}{exo:g11:subduction-zones:8}
(1) $(74.6 + 70.0 + 65.5 + 63.1 + 52.6)/5 \approx \qty{65.2}{mm}$ प्रति वर्ष;
$(47.6 + 31.0 + 29.2 + 22.8 + 19.8 + 15.2 + 14.9 + 10.6)/8 \approx
\qty{23.9}{mm}$ प्रति वर्ष; अनुपात $65.2/23.9 \approx 2.7$। (2) प्लेटें
सबसे तेज़ तब चलती हैं जब उनसे कोई \omterm{def:g11:subduction-zones:slab}{स्लैब} लटकती है: \omterm{def:g11:subduction-zones:slab-pull}{डूबती स्लैब का खिंचाव} प्लेटों को
चलाने वाला सबसे प्रबल बल है, जैसा फ़ोर्साइथ और उयेदा ने निष्कर्ष निकाला।
\end{solution}

\begin{solution}{exo:g11:subduction-zones:9}
(1) \qty{83}{mm} बराबर \qty{0.083}{m} है: प्लेट-गति के $50/0.083 \approx 600$ वर्ष और
$60/0.083 \approx 720$ वर्ष। (2) 869 से 2011 तक 1142
वर्ष हैं: उसी कोटि के। यदि महान भूकंपों के बीच प्लेट-सीमा जकड़ी रहती है,
तो 2011 की सरकन ने कई सदियों की संचित गति मुक्त की।
\end{solution}

\begin{solution}{exo:g11:subduction-zones:10}
(1) जापान: $6400 - 3200\,e^{-133/62.8} \approx \qty{6015}{m}$; टोंगा:
$6400 - 3200\,e^{-99/62.8} \approx \qty{5740}{m}$; एंडीज़:
$2500 + 350\sqrt{49} = \qty{4950}{m}$। (2) जापान तक पहुँचने वाली तली: सबसे
पुरानी, इसलिए सबसे ठंडा और सबसे घना \omterm{def:g11:lithosphere-asthenosphere:lithosphere}{स्थलमंडल}, और सबसे गहरी तली।
\end{solution}

\begin{solution}{exo:g11:subduction-zones:11}
(1) $3380 \times 1.06 \approx \qty{3583}{kg/m^3}$। (2)
$(3583 - 2900)/2900 \approx 0.24$: लगभग \qty{24}{\percent} की वृद्धि।
(3) पहले, भूपर्पटी प्रावार से
$3380 - 2900 = \qty{480}{kg/m^3}$ हल्की है; बाद में, वह लगभग
\qty{203}{kg/m^3} अधिक घनी है। वह अब \omterm{def:g11:subduction-zones:slab}{स्लैब} को ऊपर नहीं थामती: उसे नीचे खींचती है।
\end{solution}

\begin{solution}{exo:g11:subduction-zones:12}
(1) $t = 6 \times 10^6 \times 3.16 \times 10^7 \approx 1.9 \times 10^{14}$~s;
$\sqrt{10^{-6} \times 1.9 \times 10^{14}} \approx \qty{1.4e4}{m}$, लगभग
\qty{14}{km}। (2) \qty{100}{km} से बहुत कम: 6 मिलियन वर्षों में \omterm{def:g11:subduction-zones:slab}{स्लैब} की केवल
बाहरी परत गरम होती है; उसका भीतरी भाग ठंडा रहता है। (3) ठंडी चट्टान
दृढ़ होती है और टूट सकती है, जबकि उसके चारों ओर का गरम प्रावार बहता है; गहरे
भूकंप ठंडी \omterm{def:g11:subduction-zones:slab}{स्लैब} के भीतर होते हैं।
\end{solution}

\begin{solution}{exo:g11:subduction-zones:13}
(1) $\sqrt{1017^2 + 670^2} \approx \qty{1218}{km}$। (2) प्रति वर्ष \qty{64}{mm}
बराबर प्रति मिलियन वर्ष \qty{64}{km} है: $1218/64 \approx 19$ मिलियन वर्ष।
(3) टोंगा के नीचे की अपेक्षा लगभग तीन गुना अधिक: एंडीज़ के नीचे की \omterm{def:g11:subduction-zones:slab}{स्लैब}
अधिक धीरे और कम नमन पर डूबती है।
\end{solution}

\begin{solution}{exo:g11:subduction-zones:14}
(1) खाई के अनुदिश \qty{1}{m} चौड़ी एक फाँक का आयतन
$10^5 \times 10^6 \times 1 = \qty{1e11}{m^3}$ है; उसका अतिरिक्त द्रव्यमान
$36 \times 10^{11} = \qty{3.6e12}{kg}$ है। (2) उसका अतिरिक्त भार
खाई के हर मीटर के लिए $3.6 \times 10^{12} \times 9.8 \approx \qty{3.5e13}{N}$ है।
(3) \omterm{def:g11:subduction-zones:slab}{स्लैब} प्लेट के शेष भाग से जुड़ी है, जो दृढ़ है:
बल पूरी प्लेट तक पहुँचता है, जो खाई की ओर घसीटी जाती है,
और भी अधिक इसलिए कि खाइयाँ हज़ारों किलोमीटर तक चलती हैं।
\end{solution}

\begin{solution}{exo:g11:subduction-zones:15}
(1) $(23.6 - 16.5) \times 111 \approx \qty{790}{km}$। (2) \qty{790}{km} पर दर
$237 - 104 = \qty{133}{mm}$ प्रति वर्ष बढ़ती है, यानी प्रति सौ किलोमीटर प्रति वर्ष लगभग
\qty{17}{mm}। (3) अभिसरण प्रशांत प्लेट का वेग घटा चाप का वेग है।
चूँकि प्रशांत प्लेट का वेग दोनों बिंदुओं पर लगभग एक ही है, इसलिए चाप
स्वयं प्रशांत प्लेट की ओर चल रहा होगा, उत्तर की ओर अधिकाधिक तेज़ी से: उसके पीछे
का पश्च-चाप खिंच रहा है, दक्षिण की अपेक्षा उत्तर में अधिक।
\end{solution}

\begin{solution}{pb:g11:subduction-zones:1}
\textbf{1.} लगभग \qty{671}{km}।

\textbf{2.} महासागर वाली ओर प्लेट अभी डूबी नहीं है: वहाँ का \omterm{def:g11:lithosphere-asthenosphere:lithosphere}{स्थलमंडल}
सतह के पास है। गहरे भूकंप \omterm{def:g11:subduction-zones:slab}{स्लैब} के भीतर होते हैं, जो
केवल ऊपरी प्लेट के नीचे है।

\textbf{3.} $z(300) = -24 + 273 = \qty{249}{km}$ और
$z(700) = -24 + 637 = \qty{613}{km}$; $\tan\delta = 364/400 = 0.91$, इसलिए
$\delta \approx \ang{42}$।

\textbf{4.} $x = (100 + 24)/0.91 \approx \qty{136}{km}$ और
$x = (400 + 24)/0.91 \approx \qty{466}{km}$।

\textbf{5.} खाई से शुरू होने वाली इस नमन की सीधी रेखा $z = 0$ से
गुज़रती, $x = 0$ पर; यह रेखा उससे ऊपर है। \omterm{def:g11:subduction-zones:slab}{स्लैब} सीधी नहीं है: खाई के पास
उसका नमन कम है, और अधिक गहराई पर वह खड़ी होती जाती है।

\textbf{6.} प्रति वर्ष \qty{16.8}{cm}, प्रति मिलियन वर्ष \qty{168}{km}।

\textbf{7.} $168/69 \approx 2.4$ गुना तेज़।

\textbf{8.} खाई के आर-पार अभिसरण प्रशांत प्लेट का वेग घटा दूसरी ओर के
चाप का वेग है। यदि चाप प्रशांत प्लेट की ओर चलता है, जैसा वह तब करता है
जब उसके पीछे का पश्च-चाप खिंचता है, तो दोनों चालें जुड़ जाती हैं, और अभिसरण
अकेली प्रशांत प्लेट की चाल से अधिक होता है।

\textbf{9.} $6400 - 3200\,e^{-99/62.8} \approx \qty{5740}{m}$।

\textbf{10.} लगभग 60 मिलियन वर्ष बाद \omterm{def:g11:ocean-ridges:oceanic-lithosphere}{महासागरीय स्थलमंडल} समग्र रूप से अपने नीचे
के \omterm{def:g11:lithosphere-asthenosphere:lithosphere}{दुर्बलतामंडल} से अधिक घना होता है; 99 मिलियन वर्ष पर वह और भी अधिक घना
है। जैसे ही उसका किनारा नीचे जाता है, उसका भार शेष भाग को भीतर खींच लेता है।

\textbf{11.} $(670 + 24)/0.91 \approx \qty{763}{km}$।

\textbf{12.} $\sqrt{763^2 + 670^2} \approx \qty{1015}{km}$।

\textbf{13.} $1015/168 \approx 6.0$ मिलियन वर्ष।

\textbf{14.} $1015/69 \approx 14.7$ मिलियन वर्ष।

\textbf{15.} $3380 \times 1.06 \approx \qty{3583}{kg/m^3}$: भूपर्पटी, जो पहले
प्रावार से हल्की थी, उससे अधिक घनी हो जाती है, और पूरी \omterm{def:g11:subduction-zones:slab}{स्लैब} अधिक ज़ोर से
खींचती है।

\textbf{16.} $t \approx 6.0 \times 10^6 \times 3.16 \times 10^7 \approx
1.9 \times 10^{14}$~s, और $\sqrt{\kappa t} \approx \qty{14}{km}$।

\textbf{17.} \qty{100}{km} के सामने \qty{14}{km} कम है: \omterm{def:g11:subduction-zones:slab}{स्लैब} का केवल बाहरी भाग
गरम हुआ है; \qty{670}{km} तक पहुँचने पर भी उसका भीतरी भाग
ठंडा है।

\textbf{18.} ठंडी \omterm{def:g11:subduction-zones:slab}{स्लैब} अपने सिरे तक टूटने लायक दृढ़ रहती है,
जबकि उसके चारों ओर का गरम प्रावार बहता है: टोंगा के नीचे के भूकंप
\omterm{def:g11:subduction-zones:slab}{स्लैब} के भीतर, लगभग \qty{671}{km} तक, होते हैं।

\textbf{19.} भूकंपीय तरंगें प्रावार में लगभग \qty{670}{km} पर एक तीव्र
परिवर्तन दिखाती हैं: वहाँ उनकी चालें उछलकर बढ़ जाती हैं।

\textbf{20.} खाई से \qty{670}{km} तक लगभग \textbf{6 मिलियन वर्ष}।
\end{solution}
